% GENERATED FILE. Edit canonical lessons, then run npm run generate:books.
% canonical-source-hash: fnv1a64:568112217ee8e78b
% canonical-lessons: TE-C193-gita, TE-C193-vrttam, TE-C193-caturasram, TE-C193-tribhujam, TE-C193-cakram

\chapter{Line, Circle, Square, Triangle, Wheel}
\label{ch:te-line-circle-square-triangle-wheel}
\hlchaptermodality{193}

\begin{chapteropening}
\textbf{By the end of this chapter:} \emph{I can say a line, a circle, a square, a triangle and a wheel.}

Complete the last lesson of chapter 193: I can say a line, a circle, a square, a triangle and a wheel.
\end{chapteropening}

\section[\texorpdfstring{\te{గీత} (gīta)}{గీత (gīta)}]{\te{గీత} (gīta) — a line}

\label{lesson:TE-C193-gita}



\begin{quote}
Before the new one: say the Telugu for a necklace, then the Telugu for a braid.
\end{quote}

\subsection*{You'll want to know: \te{గీత}}
\textbf{\te{గీత}} — \emph{gīta} — \textquotedblleft{}a line\textquotedblright{}.

\begin{grammarlens}[title={What you've built}]
One more word for this chapter.
\end{grammarlens}

\subsection*{Your turn}
\begin{itemize}
\raggedright
  \item \emph{Say it:} \emph{gīta}
  \item \emph{Say it:} \emph{gīta}, once more
  \item \emph{Say it:} say \emph{gīta}
  \item \emph{Recall:} say the Telugu for a necklace, then the Telugu for a braid, then say \emph{gīta} again
\end{itemize}

\begin{tcolorbox}[breakable,colback=teal!4,colframe=teal!35!black,title={Before you move on}]
What does \te{గీత} mean? (\textquotedblleft{}A line\textquotedblright{}.) Say it once more.
\end{tcolorbox}
\section[\texorpdfstring{\te{వృత్తం} (vṛttaṁ)}{వృత్తం (vṛttaṁ)}]{\te{వృత్తం} (vṛttaṁ) — a circle}

\label{lesson:TE-C193-vrttam}



\begin{quote}
Before the new one: say the Telugu for a braid, then the Telugu for a line.
\end{quote}

\subsection*{You'll want to know: \te{వృత్తం}}
\textbf{\te{వృత్తం}} — \emph{vṛttaṁ} — \textquotedblleft{}a circle\textquotedblright{}.

\begin{grammarlens}[title={What you've built}]
One more word for this chapter.
\end{grammarlens}

\subsection*{Your turn}
\begin{itemize}
\raggedright
  \item \emph{Say it:} \emph{vṛttaṁ}
  \item \emph{Say it:} \emph{vṛttaṁ}, once more
  \item \emph{Say it:} say \emph{vṛttaṁ}
  \item \emph{Recall:} say the Telugu for a braid, then the Telugu for a line, then say \emph{vṛttaṁ} again
\end{itemize}

\begin{tcolorbox}[breakable,colback=teal!4,colframe=teal!35!black,title={Before you move on}]
What does \te{వృత్తం} mean? (\textquotedblleft{}A circle\textquotedblright{}.) Say it once more.
\end{tcolorbox}
\section[\texorpdfstring{\te{చతురస్రం} (caturasraṁ)}{చతురస్రం (caturasraṁ)}]{\te{చతురస్రం} (caturasraṁ) — a square}

\label{lesson:TE-C193-caturasram}



\begin{quote}
Before the new one: say the Telugu for a line, then the Telugu for a circle.
\end{quote}

\subsection*{You'll want to know: \te{చతురస్రం}}
\textbf{\te{చతురస్రం}} — \emph{caturasraṁ} — \textquotedblleft{}a square\textquotedblright{}.

\begin{grammarlens}[title={What you've built}]
One more word for this chapter.
\end{grammarlens}

\subsection*{Your turn}
\begin{itemize}
\raggedright
  \item \emph{Say it:} \emph{caturasraṁ}
  \item \emph{Say it:} \emph{caturasraṁ}, once more
  \item \emph{Say it:} say \emph{caturasraṁ}
  \item \emph{Recall:} say the Telugu for a line, then the Telugu for a circle, then say \emph{caturasraṁ} again
\end{itemize}

\begin{tcolorbox}[breakable,colback=teal!4,colframe=teal!35!black,title={Before you move on}]
What does \te{చతురస్రం} mean? (\textquotedblleft{}A square\textquotedblright{}.) Say it once more.
\end{tcolorbox}
\section[\texorpdfstring{\te{త్రిభుజం} (tribhujaṁ)}{త్రిభుజం (tribhujaṁ)}]{\te{త్రిభుజం} (tribhujaṁ) — a triangle}

\label{lesson:TE-C193-tribhujam}



\begin{quote}
Before the new one: say the Telugu for a circle, then the Telugu for a square.
\end{quote}

\subsection*{You'll want to know: \te{త్రిభుజం}}
\textbf{\te{త్రిభుజం}} — \emph{tribhujaṁ} — \textquotedblleft{}a triangle\textquotedblright{}.

\begin{grammarlens}[title={What you've built}]
One more word for this chapter.
\end{grammarlens}

\subsection*{Your turn}
\begin{itemize}
\raggedright
  \item \emph{Say it:} \emph{tribhujaṁ}
  \item \emph{Say it:} \emph{tribhujaṁ}, once more
  \item \emph{Say it:} say \emph{tribhujaṁ}
  \item \emph{Recall:} say the Telugu for a circle, then the Telugu for a square, then say \emph{tribhujaṁ} again
\end{itemize}

\begin{tcolorbox}[breakable,colback=teal!4,colframe=teal!35!black,title={Before you move on}]
What does \te{త్రిభుజం} mean? (\textquotedblleft{}A triangle\textquotedblright{}.) Say it once more.
\end{tcolorbox}
\section[\texorpdfstring{\te{చక్రం} (cakraṁ)}{చక్రం (cakraṁ)}]{\te{చక్రం} (cakraṁ) — a wheel}

\label{lesson:TE-C193-cakram}



\begin{quote}
Before the new one: say the Telugu for a square, then the Telugu for a triangle.
\end{quote}

\subsection*{You'll want to know: \te{చక్రం}}
\textbf{\te{చక్రం}} — \emph{cakraṁ} — \textquotedblleft{}a wheel\textquotedblright{}.

\begin{grammarlens}[title={What you've built}]
One more word for this chapter.
\end{grammarlens}

\subsection*{Your turn}
\begin{itemize}
\raggedright
  \item \emph{Say it:} \emph{cakraṁ}
  \item \emph{Say it:} \emph{cakraṁ}, once more
  \item \emph{Say it:} say \emph{cakraṁ}
  \item \emph{Recall:} say the Telugu for a square, then the Telugu for a triangle, then say \emph{cakraṁ} again
\end{itemize}

\begin{tcolorbox}[breakable,colback=teal!4,colframe=teal!35!black,title={Before you move on}]
What does \te{చక్రం} mean? (\textquotedblleft{}A wheel\textquotedblright{}.) Say it once more.
\end{tcolorbox}
